%%%%%%%%%%%%%%%%
% ABSTRACT
%%%%%%%%%%%%%%%%

\chapter*{Resumen}
La deriva arquitectónica es la divergencia progresiva entre la arquitectura prevista de un sistema y su implementación real, y constituye una de las principales fuentes de deuda técnica en las pequeñas y medianas empresas de software, que rara vez disponen de documentación arquitectónica actualizada o de recursos para adoptar plataformas de verificación sofisticadas. Este trabajo propone un mecanismo automatizado que detecta dicha deriva combinando el análisis estático de grafos de dependencia con la ejecución de \textit{Fitness Functions}, entendidas como verificaciones objetivas que cuantifican la conformidad estructural de un proyecto. El mecanismo construye un grafo dirigido a partir del código compilado, evalúa tres dimensiones —integridad de capas, ausencia de ciclos y control de la concentración de dependencias salientes— y calcula una Tasa de Conformidad que clasifica el estado del sistema como aceptable, en observación o no conforme, acompañando cada hallazgo de la evidencia que lo origina. La propuesta se evaluó mediante un experimento controlado sobre cinco escenarios: una implementación de referencia y cuatro versiones alteradas, tres de ellas con una única deriva inducida para medir la sensibilidad de cada regla por separado. El mecanismo detectó la totalidad de las derivas sin producir hallazgos sobre la implementación de referencia. La evaluación reveló, además, que la Tasa de Conformidad, al ser una proporción no ponderada, diluye las violaciones aisladas y puede clasificar como aceptable un sistema con una desviación real; ese hallazgo motivó un refinamiento que incorpora la gravedad del peor hallazgo junto a la extensión que mide la métrica, con el cual los cinco escenarios pasaron a clasificarse conforme a lo esperado. Los resultados evidencian la viabilidad del enfoque para la gobernanza continua sin depender de documentación arquitectónica externa, y precisan los límites que toda métrica de conformidad agregada arrastra.

{\bf Palabras Clave}: deriva arquitectónica, análisis estático, grafos de dependencia, \textit{Fitness Functions}, gobernanza arquitectónica.

\chapter*{Abstract}
Architectural drift is the progressive divergence between a system's intended architecture and its actual implementation, and it represents one of the main sources of technical debt in small and medium-sized software enterprises, which seldom keep architectural documentation up to date or can afford sophisticated verification platforms. This work proposes an automated mechanism that detects such drift by combining static dependency-graph analysis with the execution of \textit{Fitness Functions}, understood as objective checks that quantify a project's structural conformance. The mechanism builds a directed graph from compiled code, evaluates three dimensions —layer integrity, absence of cycles, and control of outgoing dependency concentration— and computes a Conformance Rate that classifies the system's state as acceptable, under observation, or non-conformant, attaching to every finding the evidence that originates it. The proposal was evaluated through a controlled experiment over five scenarios: a reference implementation and four altered versions, three of which introduce a single drift each so that the sensitivity of every rule can be measured in isolation. The mechanism detected every induced drift without reporting any finding on the reference implementation. The evaluation further revealed that the Conformance Rate, being an unweighted proportion, dilutes isolated violations and may classify as acceptable a system that contains a genuine deviation; this finding motivated a refinement that combines the severity of the worst finding with the extent that the metric measures, under which all five scenarios were classified as expected. The results evidence the feasibility of the approach for continuous governance without relying on external architectural documentation, and delimit the constraints that any aggregate conformance metric carries.

{\bf Keywords}: architectural drift, static analysis, dependency graphs, \textit{Fitness Functions}, architectural governance.
